\documentclass[10pt,a4paper]{article}
\usepackage[margin=1.8cm]{geometry}
\usepackage[T1]{fontenc}
\usepackage{lmodern,amsmath,amssymb,booktabs}
\usepackage[colorlinks=true,urlcolor=blue,citecolor=blue,linkcolor=blue]{hyperref}
\setlength{\parindent}{0pt}
\setlength{\parskip}{5pt}
\setlength{\abovedisplayskip}{6pt}
\setlength{\belowdisplayskip}{6pt}
\raggedbottom
\newcommand{\apar}{A_{\parallel}}
\newcommand{\gev}{\mathrm{GeV}}
\begin{document}
\begin{center}
{\Large Reconstruction of the HERMES Born longitudinal asymmetry plots}\par
\small 30 September 2026
\end{center}

We construct proton and deuteron projections of the published HERMES Born
longitudinal asymmetry, $\apar$, using an independent fit to unpolarized
cross sections. The original asymmetry release lacks the unpolarized cell yields
needed to determine an integrated event-sample asymmetry directly.
The eight reconstructed plots comprise both targets versus
$x$ and versus $Q^2$, with $Q^2>1$ and $Q^2>4\,\gev^2$ selections.
These are model-assisted reconstructions, conditional on the fitted weights
and the stated bin assumptions; they are not additional published HERMES measurements.

The inputs are the 45 Born $(x,Q^2)$ cells for each target in HERMES
Tables XI and XII, available as HEPData Tables 7 and 8, and the full statistical
correlation matrices in HEPData Tables 18 and 19~\cite{spin,data}.
Published Born values are used directly, without multiplying $A_1$ by an
averaged depolarization factor. Eight cells per target lie below the generator's
$Q^2=1\,\gev^2$ boundary and receive zero weight, leaving 37 supported cells.

For the Herwig cell predictions, both stored order contributions are normalized
before subtracting NEGNLO magnitudes from POSNLO in each target-component and helicity bin.
We then form
\begin{equation}
 \sigma_{\rm UU}=\frac{\sigma_{\rm PP}+\sigma_{\rm PM}+\sigma_{\rm MP}+\sigma_{\rm MM}}{4},
 \qquad
 \sigma_{\rm LL}=\frac{\sigma_{\rm PP}+\sigma_{\rm MM}-\sigma_{\rm PM}-\sigma_{\rm MP}}{4},
 \qquad \apar^{\rm MC}=\frac{\sigma_{\rm LL}}{\sigma_{\rm UU}}.
\end{equation}
Deuteron cells already use
$\sigma_{\rm UU}^d=(\sigma_{\rm UU}^p+\sigma_{\rm UU}^n)/2$ and
$\sigma_{\rm LL}^d=0.925(\sigma_{\rm LL}^p+\sigma_{\rm LL}^n)/2$.
The reconstruction does not apply this factor again. No extra depolarization,
polarization, dilution, tensor, nuclear or normalization correction is applied
to the experimental cells. The generator's existing finite-$Q^2$ limitations remain;
this procedure adds no explicit $g_2$, target-mass or twist-three model.

The HERMES GD11-P and GD11-D fits~\cite{unpolarized} supply the unpolarized
$F_2$ per nucleon. We convert it to the lepton Born cross section with Eq.~(2.6)
of Ref.~\cite{unpolarized}, using the R1998 parameterization of the virtual-photon
longitudinal-to-transverse ratio $R=\sigma_L/\sigma_T$~\cite{r1998}.
For output region $R_r$, source cell $B_i$ and the common fiducial region $\mathcal F$,
we calculate
\begin{align}
 U_{ri}&=\int_{R_r\cap B_i\cap\mathcal F} dx\,dQ^2\,
       \frac{d^2\sigma_{\rm UU}^{\rm GD11}}{dx\,dQ^2},
 & W_{ri}&=\frac{U_{ri}}{\sum_j U_{rj}},\label{eq:weights}\\
 \widehat A_r^{\rm data}&=\sum_i W_{ri}A_{\parallel,i}^{\rm data},
 &\widehat A_r^{\rm MC}&=\sum_i W_{ri}A_{\parallel,i}^{\rm MC}.
\end{align}
The same positive weights act on data and Herwig cells; no Herwig unpolarized
prediction or polarized PDF determines the experimental weights.
Thus the reconstructed Herwig curve uses the same averaging as the reconstructed
data. It differs from the original event-level Herwig ratio of integrated
cross sections, which uses Herwig's own cross-section distribution and retains
its event-level variation inside cells.

The integration uses $E_\ell=27.6\,\gev$, $1<Q^2<20\,\gev^2$,
$0.10<y<0.91$, $W^2>3.24\,\gev^2$ and
$0.04<\theta<0.22$ radians. We restrict the reconstructed integral to the
published cell union, whose supported $x$ range is $0.02124<x<0.9$.
The plot definitions are
\begin{center}
\begin{tabular}{@{}lll@{}}
\toprule
Selection & Versus $x$: integrate over $Q^2$ & Versus $Q^2$: integrate over $x$\\
\midrule
$Q^2>1\,\gev^2$ & 15 physical $x$ bins; $1<Q^2<20$ & 8 bins; $0.02124<x<0.9$\\
$Q^2>4\,\gev^2$ & Same $x$ bins; 10 supported & 4 bins; $0.02124<x<0.9$\\
\bottomrule
\end{tabular}
\end{center}
Each row is produced for both targets. The $Q^2$ edges are
$1,1.5,2,3,4,6,8,12,20\,\gev^2$; the stricter selection retains
$4,6,8,12,20\,\gev^2$. Its five lowest-$x$ bins have zero accepted fitted yield
and remain masked. The main $x$ projection merges whole source cells and needs
no interpolation of their spin asymmetries. The $Q^2$ projections and the
$Q^2>4$ $x$ projection split some cells and assume constant $\apar$ inside each
split cell, on both the experimental and Herwig sides.

Black points show $\widehat A^{\rm data}$ with statistical errors from
\begin{equation}
 C_{\rm projected}=W C_{\rm stat} W^{\mathsf T}.
\end{equation}
This retains correlations between source cells and between output bins sharing
source cells. Since only experimental systematic diagonals are public, the gray
outer bars combine the statistical error in quadrature with the conservative
bound $s_r^{\rm bound}=\sum_i|W_{ri}|s_i$. This is not a recovered systematic
covariance or a specified confidence interval. Published systematics already
contain the 5.2\% proton and 5.0\% deuteron normalization uncertainties.
Orange segments show $\widehat A^{\rm MC}$; their bands use the independent-cell
approximation $(\delta A_r^{\rm MC})^2=\sum_iW_{ri}^2(\delta A_{\parallel,i}^{\rm MC})^2$.
If any positively weighted MC cell has no finite prediction, that MC projection
is masked without changing the experimental weights or covariance.
Horizontal bars span the output bins; $x$ markers average the published cell
mean $x$ with the fitted weights, while $Q^2$ markers use geometric bin centres.

The displayed errors are conditional on central fitted weights and omit
fit-weight and within-cell spin-shape uncertainties. The public rounded GD11
parameter covariance matrices are indefinite and are not converted into a
fit-error band. The overview figures simply assemble the individual plots.
The separate sensitivity figures remove $3.24<W^2<4\,\gev^2$ from the weights,
while retaining the same source-cell asymmetries; dashed curves show the
resulting data and MC shifts. This region lies below GD11's stated fit domain
and contributes about 45\% of the central weight in the final $x$ bin.
These shifts are not one-sigma errors or measurements with a newly imposed
experimental event cut. Additional numerical controls use ALLM97 proton weights
(with GD11-D/GD11-P for the deuteron), R1990, and doubled Gauss quadrature order
from 32 to 64. They are recorded as sensitivity or convergence checks, not
added to the displayed errors.

\small
\begin{thebibliography}{4}
\setlength{\itemsep}{0pt}
\bibitem{spin} HERMES Collaboration, Phys. Rev. D \textbf{75} (2007) 012007,
\href{https://arxiv.org/abs/hep-ex/0609039}{arXiv:hep-ex/0609039}.
\bibitem{data} \href{https://www.hepdata.net/record/ins726689}{HEPData record ins726689},
Born asymmetries and statistical correlation matrices.
\bibitem{unpolarized} HERMES Collaboration, JHEP \textbf{05} (2011) 126,
\href{https://arxiv.org/abs/1103.5704}{arXiv:1103.5704}.
\bibitem{r1998} E143 Collaboration, Phys. Lett. B \textbf{452} (1999) 194--200,
\href{https://arxiv.org/abs/hep-ex/9808028}{arXiv:hep-ex/9808028}.
\end{thebibliography}
\normalsize
Full implementation details and pinned numerical inputs:
\href{https://github.com/apapaefs/HerwigPolarizedPheno/blob/main/docs/hermes-apar-integrated-data.md}{HerwigPolarizedPheno reconstruction documentation}.
\end{document}
